\documentclass[11pt,oneside]{article}
\usepackage[T1]{fontenc}
\usepackage[utf8]{inputenc}
\usepackage{lmodern,microtype}
\usepackage[a4paper,margin=27mm,headheight=15pt]{geometry}
\usepackage{amsmath,amssymb,amsthm,mathtools,bm}
\usepackage{enumitem,booktabs,longtable,array}
\usepackage[dvipsnames]{xcolor}
\usepackage{hyperref}
\usepackage{fancyhdr}
\hypersetup{colorlinks=true,linkcolor=MidnightBlue,citecolor=MidnightBlue,urlcolor=MidnightBlue,pdftitle={Surcomplex Analysis: A Hahn-Supported Theory},pdfauthor={OpenAI}}
\pagestyle{fancy}\fancyhf{}\fancyhead[L]{\small Surcomplex analysis}\fancyhead[R]{\small Hahn-supported holomorphy}\fancyfoot[C]{\thepage}
\setlength{\parskip}{4pt}\setlength{\parindent}{1.3em}
\setlist{itemsep=2pt,topsep=4pt}
\allowdisplaybreaks[2]
\newtheorem{theorem}{Theorem}[section]
\newtheorem{proposition}[theorem]{Proposition}
\newtheorem{lemma}[theorem]{Lemma}
\newtheorem{corollary}[theorem]{Corollary}
\theoremstyle{definition}\newtheorem{definition}[theorem]{Definition}
\newtheorem{example}[theorem]{Example}
\newtheorem{remark}[theorem]{Remark}
\newcommand{\No}{\mathbf{No}}
\newcommand{\Oz}{\mathbf{Oz}}
\newcommand{\K}{\mathbb K}
\newcommand{\C}{\mathbb C}
\newcommand{\R}{\mathbb R}
\newcommand{\N}{\mathbb N}
\newcommand{\Z}{\mathbb Z}
\newcommand{\HH}{\mathcal H}
\newcommand{\OO}{\mathcal O}
\newcommand{\mm}{\mathfrak m}
\newcommand{\supp}{\operatorname{supp}}
\newcommand{\st}{\operatorname{st}}
\newcommand{\ord}{\operatorname{ord}}
\newcommand{\Res}{\operatorname{Res}}
\newcommand{\ev}{\operatorname{ev}}
\newcommand{\lc}{\operatorname{lc}}
\newcommand{\id}{\operatorname{id}}
\newcommand{\dd}{\,\mathrm d}
\newcommand{\Hsum}{\mathop{\sum}\nolimits^{\mathrm H}}
\newcommand{\hatF}{\widehat F}
\newcommand{\val}{\nu}
\newcommand{\htop}{\mathrm H}
\newcommand{\Fin}{\K_{\mathrm{fin}}}
\newcommand{\Disc}{\mathbb D}
\newcommand{\abs}[1]{\left|#1\right|}
\newcommand{\norm}[1]{\left\|#1\right\|}
\newcommand{\mres}{\operatorname{res}}
\newcommand{\lead}{\operatorname{in}}
\newcommand{\Usharp}{U^{\sharp}}
\newcommand{\mc}{\mathcal}
\title{\vspace{-1cm}\textbf{Surcomplex Analysis}\\[4pt]
\Large A Hahn-Supported Theory of Holomorphic Functions on $\No(i)$\\[7pt]
\large Infinitesimal germs, Cauchy calculus, Weierstrass preparation,\\and conservation of zero multiplicities}
\author{Research article prepared with OpenAI}
\date{21 September 2026}
\begin{document}
\maketitle
\begin{abstract}
The algebraically closed class field $\K=\No(i)$ does not inherit classical complex analysis merely by replacing real coordinates with surreal coordinates. In the topology defined by all positive surreal radii, every convergent set-indexed net is eventually constant. Nonconstant bounded functions can be locally constant and hence have derivative zero everywhere. Ordinary sequential convergence and ordinary parametrized contours therefore cannot provide the desired foundations.

We develop a two-level alternative. At the microscopic level every formal power series with surcomplex coefficients has a faithful realization on a sufficiently small neighborhood, using Conway--Hahn summation. This gives local differentiation, inverse and implicit function theorems, ramification normal forms, and coordinate-invariant residues. At the macroscopic level we use well-ordered Hahn series whose coefficients are ordinary holomorphic functions on a common complex domain. Their Taylor evaluations on infinitesimal neighborhoods are compatible with genuine surreal-radius differentiation. Coefficientwise integration then yields Cauchy's theorem and formula, a residue theorem, and an argument principle.

The central structural result is a constructive Hahn--Weierstrass preparation theorem. If the leading coefficient has a zero of order $m$ at $c$, the evaluated function has exactly $m$ zeros in the entire monad of $c$, counted with multiplicity. This yields multiscale root stability and a Rouch\'e theorem. We also prove an open mapping theorem, a maximum-modulus theorem, and a sharp residue-level replacement for Liouville's theorem. The construction is a specified analytic framework, not a claim that every notion of surcomplex analyticity is equivalent. Known surreal and Hahn-theoretic foundations are explicitly distinguished from the results proved here.
\end{abstract}
\clearpage
\setcounter{tocdepth}{1}
\tableofcontents
\clearpage

\section{Scope, literature, and the meaning of the theory}\label{sec:intro}
A \emph{surcomplex number} is an element of
\[
 \K=\No(i)=\{x+iy:x,y\in\No\},\qquad i^2=-1.
\]
The motivation is to retain the algebraic and local analytic strengths of complex analysis while allowing infinitesimal, infinite, and transfinite scales. The word ``field'' below is understood in the class sense when its universe is $\K$ or $\No$. Every individual sum and every support in this article is a set.

The question is not whether $\K$ is algebraically closed: that follows from the real closedness of $\No$. The problem is to specify which functions, sums, domains, and integrals deserve to be called analytic. A definition that addresses only one of these issues can be misleading. For example, the usual difference quotient makes perfect sense using surreal radii, but it admits too many locally constant functions to enforce global uniqueness.

\subsection{What is already available}
Conway's normal form and Gonshor's development of surreal functions supply the algebraic background \cite{Conway,Gonshor}. Van den Dries and Ehrlich construct the restricted analytic structure on $\No$ and establish the elementary-extension result for restricted analytic functions and exponentiation \cite{vdDE,vdDEerr}. Thus infinitesimal evaluation of ordinary analytic germs is not new. Rubinstein-Salzedo and Swaminathan study another approach to surreal limits and analysis, using definitions different from the all-radii topology used here \cite{RSS}.

There is also an existing global surcomplex exponential. Ehrlich and Kaplan obtain a canonical exponential associated with the omnific integer part, with kernel $2\pi i\Oz$ \cite[Section 11]{EK}. Section~\ref{sec:exp} explains how its local calculus fits the present framework. Costin and Ehrlich develop extensions and integration for a broad resurgent/transserial class \cite{CE}; their global integration problem is not identical to the coefficientwise contour problem treated here.

The term \emph{Hahn-holomorphic} already occurs in work of M\"uller and Strohmaier on normally convergent generalized expansions on a logarithmic cover of the complex plane \cite{MS}. Our objects have a different direction of dependence: their Hahn monomials are fixed surreal constants, while their coefficients vary holomorphically in an ordinary complex variable. We therefore use the term \emph{Hahn-supported holomorphic function}. Hahn summability itself is standard \cite{Alling,Marker}. Ordinary complex-analytic results used below have their customary meanings and hypotheses; a reference is \cite{Ahlfors}.

\subsection{Contributions and their limits}
The purpose is to establish a coherent theory with proofs, rather than to announce that a named open problem has been settled. No priority claim is made for the entire package or for familiar formal-algebraic ingredients. The main results are the following.

\begin{center}
\begin{tabular}{>{\raggedright\arraybackslash}p{0.33\textwidth}>{\raggedright\arraybackslash}p{0.58\textwidth}}
\toprule
Result & Precise setting \\
\midrule
Local realization and inverse calculus & Arbitrary formal series over $\K$, after sufficiently small rescaling.\\
Taylor evaluation and differentiation & A common well-ordered support with holomorphic coefficient functions.\\
Weierstrass preparation and zero clusters & Infinitesimal Hahn perturbations of a complex zero of finite order.\\
Cauchy, residues, argument principle & Coefficientwise integrals on ordinary complex contours and their specified affine rescalings.\\
Open mapping and maximum modulus & Evaluated Hahn-supported functions on a connected macroscopic domain.\\
Liouville replacement & Real-bounded functions have constant residue, but arbitrary infinitesimal holomorphic corrections.\\
\bottomrule
\end{tabular}
\end{center}
There is no unrestricted ``holomorphic implies analytic'' theorem for arbitrary class functions. There is no assertion that all infinite sums are limits of their partial sums. A contour integral below is not a Riemann integral in the all-radii topology. These distinctions are hypotheses of the theory, not technical qualifications to be suppressed.

\section{Surcomplex algebra, valuation, and summation}\label{sec:foundations}
\subsection{Modulus and normal form}
Conjugation and modulus are defined by
\[
 \overline{x+iy}=x-iy,\qquad |x+iy|=\sqrt{x^2+y^2}\in\No_{\geq0}.
\]
Real closedness makes the square root unambiguous. The familiar algebraic arguments give
\[
 |zw|=|z||w|,\qquad |z+w|\leq |z|+|w|,\qquad
 z^{-1}=\bar z/|z|^2\quad(z\ne0).
\]
For example, the triangle inequality follows by squaring and applying the ordered-field Cauchy--Schwarz inequality in dimension two. No completeness is used.

Write
\[
 t^\gamma:=\omega^{-\gamma}\qquad(\gamma\in\No).
\]
Here $t^\gamma$ is notation for a Conway monomial, not a use of a newly defined complex exponential. Normal form gives a unique expression
\begin{equation}\label{eq:normal}
 z=\Hsum_{\gamma\in S} c_\gamma t^\gamma,
 \qquad c_\gamma\in\C\setminus\{0\},
\end{equation}
where $S\subseteq\No$ is a well-ordered set. Combining the real and imaginary normal forms proves the complex version. We use the standard normal-form and Hahn-field theorems in this form \cite{Conway,Gonshor,vdDE}.

For nonzero $z$, put
\[
 \val(z)=\min S,\qquad \lc(z)=c_{\val(z)},\qquad \val(0)=+\infty.
\]
Then
\[
 \val(zw)=\val(z)+\val(w),\qquad
 \val(z+w)\geq\min(\val(z),\val(w)).
\]
The relation $z\prec w$ means $|z|<|w|/n$ for every positive integer $n$. For $z,w\ne0$ it is equivalent to $\val(z)>\val(w)$. In particular, $z\prec1$ means that $z$ is infinitesimal.

\begin{definition}
The finite elements and infinitesimals are
\[
 \Fin=\{z:|z|<n\text{ for some }n\in\N\},\qquad
 \mm=\{z:|z|<1/n\text{ for every }n\geq1\}.
\]
Every finite element is uniquely $c+\epsilon$, with $c\in\C$ and $\epsilon\in\mm$. We write $\st(c+\epsilon)=c$ and call $c+\mm$ the \emph{monad} of $c$.
\end{definition}
Indeed, a finite normal form has no negative exponent. Its exponent-zero coefficient is the standard part. Equivalently, $\Fin$ is the valuation ring and $\mm$ its maximal ideal, with residue field $\C$.

\subsection{Strong summability}
\begin{definition}\label{def:sum}
A set-indexed family $(z_j)_{j\in J}$ is \emph{Hahn summable} if the union of its supports is well ordered and each exponent belongs to only finitely many of those supports. Its Hahn sum is obtained by adding coefficients at each exponent. We also say \emph{strongly summable}.
\end{definition}
Repeated summation is legitimate only after this condition, or an appropriate jointly indexed version of it, has been checked. Ordinary convergence of numerical coefficients cannot replace the finite-incidence requirement.

\begin{lemma}[Hahn--Neumann support lemma]\label{lem:neumann}
Let $S,T$ be well-ordered subsets of an ordered abelian group.
The sum $S+T$ is well ordered, and each element has only finitely many representations $s+t$ with $s\in S,t\in T$.
If $A$ is well ordered and contained in the strictly positive cone, the monoid
\[
 [A]=\{a_1+\cdots+a_n:n\geq0,\ a_j\in A\}
\]
is well ordered, and every element has only finitely many representations by finite ordered words in $A$. Consequently $S+[A]$ is well ordered with the corresponding finite-incidence property.
\end{lemma}
This standard support lemma is the combinatorial foundation of Hahn fields; see \cite[Section 7.21]{Alling} and \cite[Lemma 2.36]{Marker}. For the first assertion, an infinite descending sequence of sums can be thinned so that its $S$-coordinates are nondecreasing, forcing a descending sequence in $T$. The same argument rules out infinitely many pairs with a fixed sum. The positive-word assertion is Neumann's lemma. We use no ordinal estimate for the order type of $[A]$; in particular, the corrected bounds in \cite{vdDEerr} cause no issue here.

\begin{corollary}\label{cor:formal-eps}
For every $\epsilon\in\mm$ and every sequence $c_n\in\C$, the family $(c_n\epsilon^n)_{n\geq0}$ is Hahn summable. More generally, a finite tuple of infinitesimals can be substituted into any formal power series over $\C$.
\end{corollary}
\begin{proof}
Use $A=\supp(\epsilon)\subseteq\No_{>0}$, or the finite union of the supports of the tuple, and apply Lemma~\ref{lem:neumann}. Scalar complex coefficients add no exponents.
\end{proof}
Thus even $\Hsum n!\epsilon^n$ exists for every infinitesimal $\epsilon$. This does not make $\sum n!z^n$ a convergent complex power series on a real-radius disk.

\subsection{Set-sized working fields and support bounds}
For a set subgroup $\Gamma\subseteq\No$, the Hahn field $\C((t^\Gamma))$ embeds canonically into $\K$. Every calculation involving a set of given normal forms takes place in such a field after enlarging $\Gamma$ to contain their supports. Divisible hulls may be taken when needed. We do not assume that an arbitrary fixed set-sized field has the bounding property of the entire surreal class.

\begin{lemma}[The set-bound principle]\label{lem:setbound}
Every set of surreal numbers has a strict surreal upper bound and a strict surreal lower bound. Every set of positive surreal numbers has a positive surreal strict lower bound.
\end{lemma}
\begin{proof}
For a set $A\subseteq\No$, the cut $\{A\mid\varnothing\}$ supplies a strict upper bound, and $\{\varnothing\mid A\}$ supplies a strict lower bound. For a positive set $P$, the cut $\{0\mid P\}$ is positive and smaller than every member of $P$.
\end{proof}
Whenever support sets are combined below, they remain sets. The set-bound principle can therefore be applied to all their exponents at once. This unusually strong freedom to choose a still smaller scale is central to the local realization theorem.

\section{Why ordinary topological complex analysis does not transfer}\label{sec:topology}
Give $\K$ the class topology generated by the balls
\[
 B(a,r)=\{z:|z-a|<r\},\qquad r\in\No_{>0}.
\]
Limits mean the usual epsilon--delta statements with \emph{all} positive surreal epsilons. Valuation balls $\{z:\val(z-a)>\beta\}$ give the same local topology: each contains a sufficiently small modulus ball, and each modulus ball contains a sufficiently small valuation ball.

\begin{theorem}[Set-indexed convergence collapses]\label{thm:topology}
Every set-indexed convergent net in $\K$ is eventually constant. Every set-indexed Cauchy net is eventually constant. Every set-sized subset of $\K$ is closed and discrete.
\end{theorem}
\begin{proof}
If $z_d\to a$, the nonzero distances $|z_d-a|$ form a set of positive surreals. Choose $r>0$ smaller than all of them. Eventually $|z_d-a|<r$, so eventually $z_d=a$.
For the Cauchy assertion, apply the same argument to all nonzero pairwise distances $|z_d-z_e|$. A tail of diameter less than their positive lower bound consists of a single value.

For a set $A$, all nonzero distances between its points have a common positive strict lower bound, proving discreteness. For $x\notin A$, the positive set $\{|x-a|:a\in A\}$ has a positive strict lower bound, giving a neighborhood of $x$ disjoint from $A$. Thus $A$ is closed.
\end{proof}
In particular, $1/n$ does not tend to zero, and the partial sums of a nontrivial Hahn series do not converge in this topology. Even the standard copy of $\C$ has the discrete induced topology. A continuous map from the ordinary real interval $[0,1]$ into $\K$ is constant: its image is a set-sized discrete space, and a connected space has connected image.

\begin{example}[A bounded, everywhere differentiable counterexample]\label{ex:locconstant}
Define
\[
 L(z)=\begin{cases}0,&z\in\Fin,\\1,&z\notin\Fin.\end{cases}
\]
The finite and infinite regions are both open. Near a finite point take a finite-radius ball; near an infinite point take radius $|z|/2$. Hence $L$ is locally constant, is nonconstant globally, satisfies $L'=0$ everywhere, and is bounded by $1$. It is even locally represented by constant power series.
\end{example}
The example does not contradict the classical identity theorem on connected complex domains: the domain here has a radically different topology. It does show why ``differentiable everywhere'' or ``locally a power series'' alone cannot supply the global rigidity one expects from complex analysis.

\section{Microscopic analytic germs over the full surcomplex field}\label{sec:micro}
\subsection{Realizing arbitrary formal series}
\begin{theorem}[Local realization]\label{thm:realization}
For every formal series $A(X)=\sum_{n\geq0}a_nX^n\in\K[[X]]$, there is $r>0$ such that
\[
 A^{\htop}(h)=\Hsum_{n\geq0}a_nh^n
\]
is defined for $|h|<r$. It defines a differentiable function there after possible further shrinking, with derivative given by formal differentiation. The same holds in finitely many variables. Finite collections of formal addition, multiplication, differentiation, and composition identities (with zero constant term in the substituted series) hold as functional identities after a common shrinking.
\end{theorem}
\begin{proof}
Let $T=\bigcup_n\supp(a_n)$. This is a set, not necessarily well ordered. Choose $B>0$ with $-B<\gamma<B$ for all $\gamma\in T$, and choose $\sigma>2B$. Put $b_n=a_nt^{n\sigma}$. The support of $b_n$ lies in
\[
 (n\sigma-B,n\sigma+B).
\]
These intervals are strictly ordered by $n$. Thus $(b_n)$ is strongly summable: its support is an ordered sum of well-ordered sets, and different blocks have no common exponent.

If $h=t^\sigma u$ and $|u|<1$, write $u=c+\eta$ with $c\in\C$ and $\eta\in\mm$. Expanding the finite polynomial $(c+\eta)^n$ shows that the joint family for $\sum b_nu^n$ has support in
\[
 \bigcup_n\supp(b_n)+[\supp(\eta)].
\]
For a fixed resulting exponent, Neumann's lemma gives finitely many base exponents and positive words. Each base exponent belongs to just one $n$-block, where the binomial expansion is finite. This proves joint summability and hence realization with $r=t^\sigma$.

For finitely many variables, group monomials by total degree; there are finitely many at each degree. The same block argument applies, followed by the positive-word lemma for the infinitesimal parts of all normalized coordinates.

For recentering, write $h=h_0+k$ and expand each $(h_0+k)^n$. With $h_0/t^\sigma$ finite and $k/t^\sigma$ infinitesimal, the preceding joint-support argument is unchanged. Terms of degree zero and one in $k$ give the value and formal derivative. The remaining terms are $k^2R(k)$, with $\val(R(k))$ bounded below independently of sufficiently small $k$: after normalization the remainder has a fixed well-ordered base support plus positive-word supports. Hence $kR(k)\to0$ as $k\to0$ in the all-radii sense. Repeating proves termwise higher differentiation.

Finally, fix finitely many formal identities of the stated types. At each total degree only finitely many coefficient expressions occur. The supports of all their finite products, before collecting equal powers, form a set. Bound that entire set and repeat the degree-block construction at a still smaller scale. All expanded joint families are then summable, so finite-degree formal identities remain valid after evaluation. This also justifies composition, rather than treating an unjustified iterated sum as an associative operation.
\end{proof}

\begin{definition}
A \emph{microanalytic germ} at $a\in\K$ is the function germ obtained from a formal series in $z-a$ by Theorem~\ref{thm:realization}. Equality of function germs means agreement on some all-radii neighborhood of $a$.
\end{definition}

\begin{lemma}[Faithfulness and leading-term dominance]\label{lem:faithful}
If $A\ne0$ has formal order $m$, then on a sufficiently small punctured neighborhood,
\[
 A^{\htop}(h)=a_mh^m(1+\theta(h)),\qquad \theta(h)\prec1.
\]
In particular $A^{\htop}(h)\ne0$ there. The algebra of microanalytic germs at a fixed point is canonically $\K[[X]]$.
\end{lemma}
\begin{proof}
Factor $A=a_mX^m(1+\sum_{j\geq1}d_jX^j)$. Apply the degree-block argument to all $d_j$ and choose the new scale so that every shifted support for $j\geq1$ is strictly positive. For all sufficiently small $h$, the tail is a Hahn sum supported in the positive cone, hence is infinitesimal. This proves nonvanishing. Equality of realized germs now implies equality of formal series, and the algebra assertions follow from Theorem~\ref{thm:realization}.
\end{proof}
The theorem is specifically local. It does not give a preferred radius, a usual radius-of-convergence formula, or a uniform extension of every formal series to all of $\K$.

\subsection{Differentiation, Cauchy--Riemann equations, and primitives}
For a microanalytic function, ordinary sum, product, quotient, and chain rules follow from the realized formal identities. If $F=u+iv$ with $z=x+iy$, its differential is multiplication by $F'(z)$. Therefore
\[
 u_x=v_y,\qquad u_y=-v_x.
\]
Conversely, if a map is Fr\'echet differentiable over $\No$ and its real-linear differential satisfies these equations, that differential is multiplication by a surcomplex scalar; the complex difference quotient exists. This is a statement about the differential, not an implication from arbitrary differentiability to Hahn-supported analyticity.

A primitive of $\sum a_nX^n$ is $\sum a_nX^{n+1}/(n+1)$, which is again locally realizable. Primitives are unique up to a constant \emph{as microanalytic germs}. On arbitrary disconnected class domains, locally constant differences remain possible.

\begin{theorem}[Inverse and implicit function theorems]\label{thm:inverse}
If a microanalytic germ $F$ at $a$ has $F'(a)\ne0$, it has a microanalytic inverse between suitable open neighborhoods of $a$ and $F(a)$.
If $P(X,Y)\in\K[[X,Y]]$ satisfies $P(0,0)=0$ and $P_Y(0,0)\ne0$, there is a unique formal series $Y=g(X)$ with $g(0)=0$ and $P(X,g(X))=0$; it realizes the unique sufficiently small microanalytic solution.
\end{theorem}
\begin{proof}
After translating and dividing by the nonzero linear coefficient, the inverse problem is to invert $X+\sum_{n\geq2}b_nX^n$. The coefficient of degree $n$ in its composition with $X+\sum_{n\geq2}c_nX^n$ is $c_n$ plus an expression in earlier coefficients. Successive cancellation constructs a unique inverse. Both composition identities hold on sufficiently small neighborhoods by Theorem~\ref{thm:realization}.

For the implicit problem, the coefficient of degree $n$ in $P(X,g(X))$ is $P_Y(0,0)g_n$ plus previously determined terms. This determines $g_n$. Formal division after substituting $Y=g(X)+Z$ gives
\[
 P(X,Y)=(Y-g(X))V(X,Y),\qquad V(0,0)=P_Y(0,0)\ne0.
\]
After shrinking, $V$ is nonzero by the multivariable version of leading-term dominance for a unit. The formal identities realize, proving local existence and uniqueness.
\end{proof}
The inverse proof can be made particularly transparent on sufficiently small valuation balls: a normalized map $X+O(X^2)$ and its inverse preserve the valuation of each input, and therefore both map the same small valuation ball into itself.

\begin{theorem}[Local ramification normal form]\label{thm:ramification}
Let a nonconstant microanalytic germ at $a$ satisfy
\[
 F(a+X)-F(a)=bX^m+O(X^{m+1}),\qquad b\ne0.
\]
There is a microanalytic coordinate $\phi(X)=X+O(X^2)$ such that
\begin{equation}\label{eq:ramification}
 F(a+X)=F(a)+b\phi(X)^m.
\end{equation}
For all sufficiently large $\beta\in\No$, the map sends
\[
 \{a+X:\val(X)>\beta\}
 \quad\hbox{onto}\quad
 \{F(a)+Y:\val(Y)>\val(b)+m\beta\},
\]
and every target other than $F(a)$ has exactly $m$ preimages in that source ball, counted without multiplicity.
\end{theorem}
\begin{proof}
Write the normalized unit as $U(X)=1+O(X)$. Its unique formal $m$th root with constant term $1$ is supplied by the binomial series. Set $\phi=XU^{1/m}$ and apply the inverse theorem. Shrink so that both $\phi(X)/X-1$ and its inverse-coordinate analogue are infinitesimal for $\val(X)>\beta$. They preserve valuation and are inverse bijections of that valuation ball. Equation~\eqref{eq:ramification} reduces the claim to $Y=bW^m$. Algebraic closedness provides exactly $m$ distinct roots for $Y\ne0$, and each has valuation $(\val(Y)-\val(b))/m>\beta$.
\end{proof}
The conclusion concerns \emph{valuation} balls. It would generally be false to claim that $X+X^2$ maps a fixed modulus ball $|X|<r$ onto itself, or that $|1+\epsilon|=1$ for an infinitesimal $\epsilon$. The latter has valuation zero but need not have modulus exactly one.

\subsection{Formal residues}
A meromorphic microgerm is a Laurent series with finitely many negative powers. Define
\[
 \Res_{X=0}\left(\sum_{n\geq-N}a_nX^n\dd X\right)=a_{-1}.
\]
\begin{theorem}[Residue and coordinate invariance]\label{thm:formalres}
A meromorphic microgerm $1$-form has a meromorphic primitive if and only if its residue is zero. If $X=\phi(Y)=YV(Y)$ with $V(0)\ne0$, residue is unchanged under pullback. If $X=Y^mV(Y)$, residue is multiplied by $m$.
\end{theorem}
\begin{proof}
All powers except $X^{-1}$ integrate by division by $n+1$, while a derivative has zero $X^{-1}$ coefficient. For coordinate change it suffices to check monomials. For $n\ne-1$, $\phi^n\phi'=(\phi^{n+1})'/(n+1)$ has zero residue. For $n=-1$,
\[
 \frac{\phi'}{\phi}=\frac{m}{Y}+\frac{V'}{V}.
\]
The second term is a power series. Only finitely many negative source powers exist, so this computation extends to every Laurent series without an illicit infinite rearrangement.
\end{proof}
Thus the quotient of meromorphic microgerm $1$-forms by exact forms is a one-dimensional $\K$-vector space, detected by residue. This is an algebraic analogue of the local obstruction represented classically by $\dd z/z$.

\section{Hahn-supported holomorphic functions}\label{sec:H}
\subsection{Definition and evaluation on monads}
Let $U\subseteq\C$ be a nonempty connected open set. Define
\[
 U^\sharp=\{z\in\Fin:\st(z)\in U\}.
\]
This is an open subclass of $\K$, but is not connected in the all-radii topology. Its \emph{macroscopic domain} is the connected complex set $U$.

\begin{definition}\label{def:H}
The algebra $\HH(U)$ consists of coefficient families
\begin{equation}\label{eq:H}
 F=\Hsum_{\gamma\in S}f_\gamma\,t^\gamma,
 \qquad f_\gamma\in\OO(U),
\end{equation}
with a common well-ordered set support $S$. The zero coefficient functions are omitted. The monomials are constants with respect to the variable in $U$.
Its valuation and leading coefficient are
\[
 \val(F)=\min S,\qquad \lead(F)=f_{\val(F)}.
\]
Write $\HH_{>0}(U)$ for zero together with the elements whose supports are strictly positive, and similarly define $\HH_{\geq0}(U)$.
\end{definition}
The common support is an essential part of the definition. A pointwise collection of unrelated surcomplex values is not automatically in $\HH(U)$. Addition and multiplication are Hahn addition and convolution with coefficients in $\OO(U)$. Connectedness makes $\OO(U)$, hence $\HH(U)$, an integral domain.

\begin{theorem}[Canonical Taylor evaluation]\label{thm:evaluation}
For $F\in\HH(U)$ and $z=c+\epsilon\in U^\sharp$, the joint family in
\begin{equation}\label{eq:evaluation}
 \widehat F(c+\epsilon)
 =\Hsum_{\gamma\in S,\ n\geq0}
   \frac{f_\gamma^{(n)}(c)}{n!}\,t^\gamma\epsilon^n
\end{equation}
is summable. This defines an injective $\K$-algebra homomorphism from $\HH(U)$ to class functions on $U^\sharp$. It commutes with restriction and conjugation of coefficients and variables.
\end{theorem}
\begin{proof}
If $\epsilon=0$, this is just a normal form with some coefficients possibly zero. Otherwise put $A=\supp(\epsilon)\subseteq\No_{>0}$. The expanded support is contained in $S+[A]$. Lemma~\ref{lem:neumann} gives both well-ordering and finite incidence, including the index $n$ because a positive word has finite length and finitely many representations. This proves existence of the joint sum.

Taylor expansion is multiplicative over $\C$, and the same support argument justifies the joint convolution for products. Constant coefficient functions evaluate to the original elements of $\K$. These facts prove the algebra assertion. If $\hatF=0$, its values at every ordinary $c\in U$ have normal form $\sum f_\gamma(c)t^\gamma=0$. Uniqueness of normal form implies $f_\gamma(c)=0$ for all $c$, hence $F=0$. Restriction and conjugation are immediate from the formula.
\end{proof}
For an ordinary holomorphic $f$, write $f^\sharp$ for the evaluation of the one-term Hahn function $f$. This is the usual infinitesimal Taylor extension of $f$.

\begin{proposition}[Compatibility with ordinary analytic operations]\label{prop:lift}
For ordinary holomorphic functions on appropriate domains,
\[
 (f+g)^\sharp=f^\sharp+g^\sharp,\quad
 (fg)^\sharp=f^\sharp g^\sharp,\quad
 (f\circ g)^\sharp=f^\sharp\circ g^\sharp.
\]
If $f$ has no zeros, $(1/f)^\sharp=1/f^\sharp$.
\end{proposition}
\begin{proof}
The first two statements follow from Theorem~\ref{thm:evaluation}. For composition, write $g^\sharp(c+\epsilon)=g(c)+\delta$ with $\delta\in\mm$, and use the formal Taylor composition identity. All terms are supported in the positive monoid generated by $\supp(\epsilon)$; finite incidence permits reassociation. The reciprocal follows either from the product identity or from the ordinary analytic reciprocal germ.
\end{proof}

\subsection{Actual differentiation, not just a formal symbol}
Define
\[
 DF=\Hsum_\gamma f_\gamma' t^\gamma.
\]
\begin{theorem}[Taylor and derivative compatibility]\label{thm:derivative}
Every evaluated $F\in\HH(U)$ is microanalytic at every $z\in U^\sharp$. For infinitesimal $h$ one has the strongly summable identity
\begin{equation}\label{eq:HTaylor}
 \hatF(z+h)=\Hsum_{k\geq0}\frac{\widehat{D^kF}(z)}{k!}h^k.
\end{equation}
Its all-radii derivative is $\hatF'(z)=\widehat{DF}(z)$, and the same holds for every higher derivative.
\end{theorem}
\begin{proof}
Write $z=c+\epsilon$. Expand $f_\gamma(c+\epsilon+h)$ formally in the two infinitesimals $\epsilon,h$. Its joint support is contained in
\[
 S+[\supp(\epsilon)\cup\supp(h)].
\]
This is a positive-word construction, so rearranging by the power of $h$ is legitimate and gives \eqref{eq:HTaylor}. It supplies the microanalytic germ.

Let $\alpha=\min S$; the zero function is trivial. After subtracting the constant and linear terms and dividing by $h$, the remaining joint family has valuation at least
\[
 \alpha+\val(h).
\]
Indeed, every remaining term contains at least one factor $h$ and any factors $\epsilon$ only increase valuation. Given an error tolerance $\eta>0$, choose a surreal $\beta>0$ with $\alpha+\beta>\val(\eta)$. For $0<|h|<t^\beta$, the error has valuation greater than $\val(\eta)$ and hence modulus less than $\eta$. This is precisely the all-radii derivative. Apply the same argument to $D^kF$.
\end{proof}
There is no conflict with Theorem~\ref{thm:topology}: $h\to0$ is a class-neighborhood limit, not the limit of a cofinal sequence of small increments.

\begin{theorem}[Macroscopic identity theorem]\label{thm:identity}
If $F,G\in\HH(U)$ agree at ordinary points in a subset of $U$ having a complex accumulation point in $U$, then $F=G$ and $\hatF=\widehat G$ on $U^\sharp$.
\end{theorem}
\begin{proof}
At ordinary points, equality of evaluated values is equality of normal forms. Every coefficient function of $F-G$ vanishes on the given set. The ordinary complex identity theorem makes every coefficient identically zero.
\end{proof}
The accumulation point in this statement is in the \emph{ordinary complex topology}, not the all-radii topology.

\subsection{Positive tails, units, and logarithms}
\begin{proposition}\label{prop:positive}
If $R\in\HH_{>0}(U)$, then $\widehat R(z)\in\mm$ for every $z\in U^\sharp$, and
\[
 (1+R)^{-1}=\Hsum_{n\geq0}(-R)^n,\qquad
 \log(1+R)=\Hsum_{n\geq1}\frac{(-1)^{n+1}}n R^n
\]
are well-defined in $\HH(U)$. Moreover
\[
 D\log(1+R)=\frac{DR}{1+R}.
\]
If $F=t^\alpha f(1+R)$ with $f$ nowhere zero in $U$, then $F$ is a unit of $\HH(U)$ and its evaluated inverse is $1/\hatF$.
\end{proposition}
\begin{proof}
All supports used for evaluation remain positive. If $A=\supp(R)$, the displayed power series have support in $[A]$ and satisfy finite incidence by Neumann's lemma. Their formal identities hold coefficientwise. Differentiation acts on the coefficient functions, not on the exponents; it does not enlarge support. The inverse formula is therefore compatible with evaluation.
\end{proof}

\section{Hahn--Weierstrass preparation and division}\label{sec:weierstrass}
This section supplies the mechanism behind actual surcomplex zero counting. It is important that no rank-one completeness argument, and no claim that $n\alpha$ is cofinal in $\No$, is used.

\subsection{Preparation by positive-degree recursion}
Let $D=\{X\in\C:|X|<r\}$, and let $m\geq1$. Define continuous complex-linear operators on $\OO(D)$ by
\[
 J_mh=\sum_{j=0}^{m-1}\frac{h^{(j)}(0)}{j!}X^j,
 \qquad
 Q_mh=\frac{h-J_mh}{X^m}.
\]
The second expression has a removable singularity at zero and is holomorphic on the same disk. Extend both operators coefficientwise to $\HH(D)$. They never introduce new Hahn exponents.

\begin{theorem}[Constructive Hahn--Weierstrass preparation]\label{thm:preparation}
For every $E\in\HH_{>0}(D)$ there are unique
\[
 p(X)=\sum_{j=0}^{m-1}p_jX^j\in\mm[X],
 \qquad V=1+v,\quad v\in\HH_{>0}(D),
\]
such that
\begin{equation}\label{eq:prep}
 X^m+E(X)=(X^m+p(X))V(X).
\end{equation}
Their supports are contained in the positive monoid generated by $\supp(E)$, apart from the unit coefficient at exponent zero.
\end{theorem}
\begin{proof}
The equation is equivalent to
\[
 p=J_m(E-pv),\qquad v=Q_m(E-pv).
\]
Introduce a bookkeeping degree that counts occurrences of $E$, not a limit parameter. Set
\[
 p^{[1]}=J_mE,\qquad v^{[1]}=Q_mE,
\]
and for $n\geq2$ set
\begin{align*}
 p^{[n]}&=-J_m\sum_{j=1}^{n-1}p^{[j]}v^{[n-j]},\\
 v^{[n]}&=-Q_m\sum_{j=1}^{n-1}p^{[j]}v^{[n-j]}.
\end{align*}
If $A=\supp(E)$, both degree-$n$ terms have support in $nA$. The family over all $n$ is jointly summable: a fixed exponent has only finitely many positive-word representations in $A$, hence only finitely many possible degrees and coefficient products. Thus
\[
 p=\Hsum_{n\geq1}p^{[n]},\qquad v=\Hsum_{n\geq1}v^{[n]}
\]
exist. Every coefficient of $p$ is a polynomial of degree less than $m$, so $p\in\mm[X]$; every coefficient of $v$ is holomorphic on $D$. Reassociation is justified by the same positive-word condition and gives \eqref{eq:prep}.

For uniqueness, suppose $(p,v)$ and $(\widetilde p,\widetilde v)$ are two positive solutions. If they differ, let
\[
 \delta=\min\bigl(\val(p-\widetilde p),\val(v-\widetilde v)\bigr).
\]
Subtracting the fixed-point equations expresses both differences as $J_m$ or $Q_m$ applied to
\[
 -(p-\widetilde p)v-\widetilde p(v-\widetilde v).
\]
Every product has valuation strictly greater than $\delta$, since $v$ and $\widetilde p$ have positive valuation. The operators do not decrease valuation. This contradicts the definition of $\delta$.
\end{proof}

\begin{remark}
The recursion is effective at each Hahn exponent: only finitely many coefficient operations contribute there. It is not necessary that the sequence of successive approximations converge in the topology of Section~\ref{sec:topology}. In fact, except in degenerate cases it does not.
\end{remark}

\subsection{Division by the prepared polynomial}
\begin{theorem}[Hahn--Weierstrass division]\label{thm:division}
Let $P=X^m+p$ with $p\in\mm[X]$ and $\deg p<m$. Every $N\in\HH(D)$ has a unique decomposition
\[
 N=PQ+R,\qquad Q\in\HH(D),\quad R\in\K[X],\quad\deg R<m.
\]
\end{theorem}
\begin{proof}
The equation for the quotient is
\[
 Q=Q_mN-Q_m(pQ).
\]
Let $T(Q)=-Q_m(pQ)$ and put
\[
 Q=\Hsum_{n\geq0}T^n(Q_mN).
\]
If $S=\supp(N)$ and $A=\supp(p)\subseteq\No_{>0}$, the terms have support in $S+nA$. Lemma~\ref{lem:neumann} proves joint summability. The quotient equation holds, and
\[
 R=J_m(N-pQ)
\]
has the required degree. If $P\Delta Q+\Delta R=0$, then
$\Delta Q=-Q_m(p\Delta Q)$, impossible for nonzero $\Delta Q$ because the right side has strictly higher valuation. Thus $\Delta Q=0$ and $\Delta R=0$.
\end{proof}

\subsection{The zero-cluster theorem}
\begin{theorem}[Exact conservation of zero multiplicity]\label{thm:cluster}
Let $0\ne F\in\HH(U)$, put $\alpha=\val(F)$ and $f=\lead(F)$, and fix $c\in U$.
If $f(c)\ne0$, $\hatF$ has no zero in $c+\mm$.
If $f$ has a zero of order $m$ at $c$, then $\hatF$ has exactly $m$ zeros in the entire monad $c+\mm$, counted with microanalytic multiplicity. More precisely, on a complex disk $D$ about $c$,
\begin{equation}\label{eq:clusterfactor}
 F(z)=t^\alpha u(z)P(z-c)V(z-c),
\end{equation}
where $u\in\OO(D)$ is nowhere zero, $V=1+v$ has positive Hahn tail, and
\[
 P(X)=X^m+p_{m-1}X^{m-1}+\cdots+p_0,\qquad p_j\in\mm.
\]
All roots of $P$ are infinitesimal.
\end{theorem}
\begin{proof}
If $f(c)\ne0$, Taylor evaluation has leading term $t^\alpha f(c)$, since all contributions involving the displacement or a higher Hahn exponent have higher valuation. Thus no cancellation can make it zero.

Otherwise choose a disk on which $f(c+X)=X^mu(c+X)$ with $u$ nowhere zero. Write
\[
 \frac{t^{-\alpha}F(c+X)}{u(c+X)}=X^m+E(X),\qquad E\in\HH_{>0}(D-c),
\]
and apply Theorem~\ref{thm:preparation}. Both $u^\sharp$ and $\widehat V$ are nonzero throughout the monad: the first has nonzero standard part and the second is $1$ plus an infinitesimal.

Since $\K$ is algebraically closed, $P$ has $m$ roots counted with multiplicity. A root $a$ with $\val(a)\leq0$ is impossible. The monic term has valuation $m\val(a)$, while every nonzero lower term satisfies
\[
 \val(p_ja^j)=\val(p_j)+j\val(a)>m\val(a).
\]
It therefore cannot cancel the leading term. All roots lie in $\mm$. Evaluation of \eqref{eq:clusterfactor} shows that these are exactly the zeros in $c+\mm$. The nonzero factors are microanalytic units at each root, so their multiplication leaves multiplicity unchanged.
\end{proof}

\begin{corollary}[Simple-zero lifting and units]\label{cor:simple}
A simple zero of the leading coefficient has a unique surcomplex zero in its monad. An element of $\HH(U)$ is a unit if and only if its leading coefficient has no zeros in $U$, equivalently if and only if its evaluation has no zeros in $U^\sharp$.
\end{corollary}
\begin{proof}
The first statement is the case $m=1$. If the leading coefficient is zero-free, Proposition~\ref{prop:positive} gives an inverse. If it has a zero, the cluster theorem produces an actual zero of the evaluated function, which is impossible for a unit.
\end{proof}

\begin{corollary}[A stronger identity principle]\label{cor:localidentity}
If an evaluated $F\in\HH(U)$ is constant on any nonempty all-radii open ball contained in $U^\sharp$, then $F$ is a constant element of $\K$ throughout $U^\sharp$.
\end{corollary}
\begin{proof}
Subtract the constant, which is itself in $\HH(U)$. If the result were nonzero, the cluster theorem would give only finitely many zeros in each monad. A sufficiently small subball of the given ball lies in one monad and contains infinitely many distinct points, a contradiction.
\end{proof}
Thus macroscopic coefficient holomorphy restores a form of unique continuation which arbitrary locally power-series functions do not possess.

\section{Coefficientwise contour integration and Cauchy's theorems}\label{sec:cauchy}
\subsection{The integral and what it does not mean}
Let $\Gamma$ be an ordinary piecewise $C^1$ complex curve and suppose
\[
 A(\zeta)=\Hsum_{\gamma\in S}a_\gamma(\zeta)t^\gamma
\]
has a common well-ordered support and piecewise continuous coefficients along $\Gamma$. Define
\begin{equation}\label{eq:integral}
 \int_\Gamma^{\htop}A(\zeta)\dd\zeta
 :=\Hsum_{\gamma\in S}
 \left(\int_\Gamma a_\gamma(\zeta)\dd\zeta\right)t^\gamma.
\end{equation}
The sum exists because its support is a subset of $S$. It is independent of the chosen ordinary piecewise $C^1$ parametrization and is $\K$-linear. To check linearity for a general surcomplex scalar, use the finite convolution property for its support and $S$.

The superscript $\htop$ will be retained to avoid suggesting a topological Riemann integral. The standard circle is not a continuous path in the all-radii topology. Formula~\eqref{eq:integral} is an additional analytic operation specified by coefficient data, and not a claim to the contrary.

\begin{theorem}[Cauchy's theorem and primitives]\label{thm:cauchyth}
If $F\in\HH(U)$ and $\Gamma$ is an ordinary closed curve null-homologous in $U$, then
\[
 \int_\Gamma^{\htop}F(\zeta)\dd\zeta=0.
\]
If $U$ is simply connected, $F$ has a primitive $G\in\HH(U)$, and for ordinary endpoints $a,b$,
\[
 \int_a^b{}^{\htop}F(\zeta)\dd\zeta=G(b)-G(a).
\]
Such a primitive is unique up to a constant in $\K$.
\end{theorem}
\begin{proof}
Apply the ordinary Cauchy theorem to each coefficient. On a simply connected domain choose the primitive of each $f_\gamma$ to vanish at a fixed ordinary basepoint. Their Hahn sum has the same support and gives $G$. If $DG=0$, every coefficient of $G$ is a constant complex function, so $G$ is a constant element of $\K$.
\end{proof}
The uniqueness statement uses coefficient holomorphy on the common connected domain. It is stronger than uniqueness for arbitrary differentiable functions on $U^\sharp$.

\subsection{Cauchy's formula at a genuinely surcomplex point}
Let $\Omega\Subset U$ be a bounded complex Jordan domain with piecewise $C^1$ positively oriented boundary $\Gamma$.

\begin{theorem}[Surcomplex Cauchy formula]\label{thm:cauchyformula}
For $F\in\HH(U)$, $z\in\Omega^\sharp$, and $k\geq0$,
\begin{equation}\label{eq:cauchyformula}
 \widehat{D^kF}(z)
 =\frac{k!}{2\pi i}\int_\Gamma^{\htop}
     \frac{F(\zeta)}{(\zeta-z)^{k+1}}\dd\zeta.
\end{equation}
The kernel on the right has a common Hahn-supported coefficient expansion in a neighborhood of $\Gamma$, so the integral is well defined.
\end{theorem}
\begin{proof}
Write $z=c+\epsilon$, with $c\in\Omega$ and $\epsilon\in\mm$. On a neighborhood of $\Gamma$,
\[
 \frac1{(\zeta-c-\epsilon)^{k+1}}
 =\Hsum_{n\geq0}\binom{n+k}{k}
       \frac{\epsilon^n}{(\zeta-c)^{n+k+1}}.
\]
The joint expansion after multiplication by $F$ is supported in $S+[\supp(\epsilon)]$, with finite incidence. Ordinary Cauchy differentiation gives
\[
 \frac{k!}{2\pi i}\int_\Gamma
  \binom{n+k}{k}\frac{f_\gamma(\zeta)}{(\zeta-c)^{n+k+1}}\dd\zeta
 =\frac{f_\gamma^{(n+k)}(c)}{n!}.
\]
At every Hahn exponent only finitely many terms are being integrated, so integration and the joint sum interchange by definition. The result is exactly the Taylor evaluation of $D^kF$.
\end{proof}
The target need not be an ordinary point. For example, both the pole position and the value extracted by the kernel may have arbitrarily many surreal scales.

\subsection{Cauchy estimates with the appropriate bound}
For a circle $|\zeta-c|=r$ contained with its interior in $U$, put
\[
 M_{\gamma,r}=\max_{|\zeta-c|=r}|f_\gamma(\zeta)|,
 \qquad
 \mathcal M_r(F)=\Hsum_\gamma M_{\gamma,r}t^\gamma\in\No_{\geq0}.
\]
\begin{proposition}[Coefficient-majorant Cauchy estimate]\label{prop:estimate}
At an ordinary center $c$,
\[
 \left|\frac{\widehat{D^nF}(c)}{n!}\right|
 \leq\frac{\mathcal M_r(F)}{r^n}.
\]
More precisely, the complex modulus of every coefficient of the left-hand Taylor coefficient is bounded by the corresponding coefficient of the right-hand majorant.
\end{proposition}
\begin{proof}
The coefficient assertion is the ordinary Cauchy estimate. For any Hahn number $a=\sum a_\gamma t^\gamma$, its coefficient absolute majorant $B=\sum|a_\gamma|t^\gamma$ satisfies $|a|\leq B$. Indeed, every coefficient of
\[
 B^2-a\bar a
 =\Hsum_{\gamma,\delta}
 \bigl(|a_\gamma||a_\delta|-\Re(a_\gamma\bar a_\delta)\bigr)t^{\gamma+\delta}
\]
is nonnegative. Thus $B^2\geq|a|^2$ in the ordered field and $B\geq|a|$. Increasing its individual nonnegative coefficients preserves the order. Apply this to the Taylor coefficient.
\end{proof}
An order bound on the \emph{whole} function can hide every subleading coefficient. Consequently the ordinary passage from Cauchy estimates to a global Liouville theorem must be reconsidered; Section~\ref{sec:liouville} does this exactly.

\begin{theorem}[Coefficientwise Morera theorem]\label{thm:morera}
Suppose $F=\sum f_\gamma t^\gamma$ has common well-ordered support on $U$, with every $f_\gamma$ continuous. If its Hahn integral around every ordinary triangle whose interior lies in $U$ is zero, then $F\in\HH(U)$.
\end{theorem}
\begin{proof}
Uniqueness of normal form makes every coefficient integral around every such triangle zero. On a disk, the usual triangular-path argument constructs a primitive of each continuous coefficient, and differentiating that primitive proves the coefficient holomorphic. Equivalently, apply the ordinary Morera theorem coefficientwise.
\end{proof}
This theorem supplies a nontrivial converse characterization, but only after the common support and coefficient continuity hypotheses have been imposed.

\section{Meromorphic functions and residues at surcomplex poles}\label{sec:residue}
\subsection{Fraction-meromorphic functions}
On a connected $U$, call an element of $\operatorname{Frac}\HH(U)$ \emph{fraction-meromorphic}. It is represented by $A/B$ with $A,B\in\HH(U)$ and $B\ne0$, and evaluated where $\widehat B\ne0$. Equivalent fractions give the same values. Locally at every point they are quotients of microanalytic germs, so isolated poles and their residues have the meaning of Section~\ref{sec:micro}.

Let $b=\lead(B)$. Away from the ordinary zero set of $b$, a geometric inverse of its positive tail expands $A/B$ as a Hahn series of holomorphic coefficient functions. These coefficients extend meromorphically across zeros of $b$; their pole orders need not be uniformly bounded. There is nevertheless a common well-ordered Hahn support. This is enough for coefficientwise contour deformation.

\begin{lemma}[Local partial fractions for an entire monad]\label{lem:partial}
Near an ordinary zero $c$ of $b$, choose a small complex disk $D$ as in the preparation theorem. Then
\[
 \frac{A}{B}=Q+\frac{R(X)}{P(X)},\qquad X=z-c,
\]
where $Q\in\HH(D)$, $R\in\K[X]$ has degree less than $\deg P$, and all roots of $P$ are infinitesimal. Consequently it is a sum of an evaluated Hahn-supported holomorphic function and finitely many terms
\[
 \frac{d_{a,j}}{(z-c-a)^j},\qquad a\in\mm.
\]
\end{lemma}
\begin{proof}
Write $B=t^\beta uPV$ by preparation, with $uV$ a unit of $\HH(D)$. Then $A/B=N/P$ with $N=t^{-\beta}A/(uV)\in\HH(D)$. The division theorem gives $N=PQ+R$. Factor $P$ over the algebraically closed field $\K$ and perform finite algebraic partial fractions. All its roots are infinitesimal by Theorem~\ref{thm:cluster}.
\end{proof}
The lemma is stronger than a statement about formal residues at an ordinary point: it resolves a whole monad into its actual surcomplex poles.

\begin{theorem}[Residue theorem for surcomplex poles]\label{thm:residuetheorem}
Let $A/B\in\operatorname{Frac}\HH(U)$, and let $\Omega\Subset U$ have piecewise $C^1$ Jordan boundary $\Gamma$. Assume that $\lead(B)$ is nonzero on $\Gamma$. Then the evaluated fraction has only finitely many poles in $\Omega^\sharp$, and
\begin{equation}\label{eq:residue}
 \int_\Gamma^{\htop}\frac{A(z)}{B(z)}\dd z
 =2\pi i\sum_{a\in\Omega^\sharp\text{ a pole}}
       \Res_{z=a}\left(\frac{\widehat A(z)}{\widehat B(z)}\dd z\right).
\end{equation}
The sum on the right is a finite ordinary sum of surcomplex numbers.
\end{theorem}
\begin{proof}
The leading denominator has finitely many zeros in $\overline\Omega$, none on its boundary. Every pole lies in one of their monads, and preparation bounds their number in each monad. Choose disjoint small standard circles around these ordinary zeros. The Hahn expansion of $A/B$ has coefficient functions meromorphic with singularities only at the leading-denominator zeros. Ordinary contour deformation applied coefficientwise reduces the outer integral to the sum of the small-circle integrals.

On each small circle apply Lemma~\ref{lem:partial}. The holomorphic term integrates to zero. If $X=z-c$ and $a\in\mm$, then on a standard circle about $X=0$,
\[
 \frac1{(X-a)^j}
 =\Hsum_{n\geq0}\binom{n+j-1}{j-1}\frac{a^n}{X^{n+j}}.
\]
Its integral is $2\pi i$ when $j=1$ and zero when $j>1$. These are exactly its residues at $X=a$. Multiplication by $d_{a,j}\in\K$ and the finite partial-fraction sum preserve the integral. Summing proves \eqref{eq:residue}.
\end{proof}

\begin{example}[A moving pole]
Let $\epsilon\in\mm$ and take a positively oriented standard circle around zero. Then
\[
 \frac1{z-\epsilon}=\Hsum_{n\geq0}\frac{\epsilon^n}{z^{n+1}},
 \qquad
 \int^{\htop}\frac{\dd z}{z-\epsilon}=2\pi i.
\]
Each coefficient of the expanded integrand is singular at the ordinary point $0$, but the actual evaluated function has its pole at $\epsilon$. The residue theorem identifies the aggregate coefficient residue with the residue at the displaced surcomplex pole. Conflating these two locations would obscure the geometry.
\end{example}

\subsection{Argument principle and leading-order Rouch\'e}
\begin{theorem}[Argument principle]\label{thm:argument}
Let $0\ne F\in\HH(U)$ and suppose its leading coefficient $f$ is nonzero on $\Gamma=\partial\Omega$ as above. Then
\begin{equation}\label{eq:argument}
 \frac1{2\pi i}\int_\Gamma^{\htop}\frac{DF}{F}\dd z
 =\sum_{a\in\Omega^\sharp:\hatF(a)=0}\ord_a(\hatF)
 =\sum_{c\in\Omega:f(c)=0}\ord_c(f).
\end{equation}
For a nonzero fraction-meromorphic function, the analogous integral equals the number of zeros minus the number of poles, counted with multiplicity, provided both leading numerator and leading denominator are nonzero on the boundary.
\end{theorem}
\begin{proof}
Near a zero of order $m$, the logarithmic derivative is $m/(z-a)$ plus a microanalytic function, so its residue is $m$. The residue theorem gives the first equality. The cluster theorem gives the second.

There is a useful independent coefficient calculation. On a neighborhood of $\Gamma$, write $F=t^\alpha f(1+E)$ with $E$ of positive support. Then
\[
 \frac{DF}{F}=\frac{f'}f+D\log(1+E).
\]
The logarithm is a single Hahn-supported function on that neighborhood. The integral of its derivative vanishes coefficientwise. Only the ordinary integer winding number of $f(\Gamma)$ remains. For fractions, subtract the logarithmic derivatives of numerator and denominator; canceled zeros cancel their contributions as well.
\end{proof}
All infinitesimal coefficients of this contour integral cancel exactly. The answer is an ordinary integer, although the zeros it counts need not be ordinary points.

\begin{theorem}[Leading-order Rouch\'e theorem]\label{thm:rouche}
Suppose on a neighborhood of $\overline\Omega$,
\[
 F=t^\alpha(f+R),\qquad G=t^\alpha(g+S),
 \qquad R,S\in\HH_{>0},\quad f,g\in\OO,
\]
and $|g|<|f|$ on $\Gamma=\partial\Omega$. Then $\hatF$ and $\widehat{F+G}$ have the same number of zeros in $\Omega^\sharp$, counted with multiplicity. In particular, a perturbation of strictly higher Hahn valuation preserves every individual monadic zero count.
\end{theorem}
\begin{proof}
For $0\leq s\leq1$, $f+sg$ has no boundary zeros. The ordinary argument integral of $f+sg$ is independent of $s$, since its derivative with respect to $s$ is
\[
 \frac1{2\pi i}\int_\Gamma
  \frac{\dd}{\dd z}\left(\frac{g}{f+sg}\right)\dd z=0.
\]
Thus $f$ and $f+g$ have equal ordinary zero counts. The cluster theorem transfers these counts exactly to $F$ and $F+G$. If $\val(G)>\val(F)$, the leading coefficient is unchanged, so its order at each ordinary zero is unchanged.
\end{proof}

\section{Open mapping, maximum modulus, and the exact Liouville defect}\label{sec:geometry}
\subsection{Open mapping and maximum modulus}
\begin{theorem}[Open mapping]\label{thm:open}
If $F\in\HH(U)$ is not a constant element of $\K$, then $\hatF:U^\sharp\to\K$ maps all-radii open subclasses to open subclasses.
\end{theorem}
\begin{proof}
By Corollary~\ref{cor:localidentity}, the microanalytic germ of $\hatF$ at every point is nonconstant. Every source neighborhood contains a sufficiently small valuation ball to which the local ramification theorem applies. Its image contains a valuation ball about the image point. Taking the union over source points proves openness.
\end{proof}

\begin{theorem}[Maximum-modulus principle]\label{thm:maxmod}
For nonconstant $F\in\HH(U)$, $|\hatF|$ has no local maximum at any point of $U^\sharp$ in the all-radii topology. If a local maximum occurs, $F$ is a constant element of $\K$.
\end{theorem}
\begin{proof}
Suppose the germ at $a$ is nonconstant and write
\[
 \hatF(a+h)=c+bh^m(1+\theta(h)),\qquad b\ne0,
 \quad\theta(h)\prec1
\]
for sufficiently small $h$. If $c=0$, every sufficiently small nonzero $h$ gives a larger modulus. If $c\ne0$, choose a positive infinitesimal $s$ as small as necessary, and choose an $m$th root $h$ of $cs/b$. Algebraic closedness supplies it, and increasing $\val(s)$ makes $h$ lie in any prescribed source neighborhood. Then
\[
 \hatF(a+h)=c\bigl(1+s(1+\theta(h))\bigr).
\]
Consequently
\[
 \frac{|\hatF(a+h)|^2-|c|^2}{|c|^2}
 =2s(1+\Re\theta(h))+s^2|1+\theta(h)|^2>0.
\]
Thus a local maximum forces the germ to be constant. Corollary~\ref{cor:localidentity} forces global constancy in $\HH(U)$.
\end{proof}
The maximum-modulus theorem does not require the image modulus to attain a supremum on a compact surcomplex set. It is a local algebraic perturbation argument, and survives even though compactness-based classical proofs are unavailable.

\subsection{A sharp replacement for Liouville}\label{sec:liouville}
Call $F\in\HH(\C)$ \emph{macroscopically entire}. This means that its coefficient functions are entire; its evaluated domain here is $\Fin$, not automatically all of $\K$.

\begin{theorem}[Residue-level Liouville theorem]\label{thm:liouville}
For $F\in\HH(\C)$, the following are equivalent:
\begin{enumerate}[label=(\roman*)]
\item There is a positive real $M$ with $|\hatF(z)|\leq M$ for every $z\in\Fin$.
\item There are $c\in\C$ and $R\in\HH_{>0}(\C)$ such that $F=c+R$.
\end{enumerate}
Thus the algebra of real-bounded macroscopically entire functions is exactly
\begin{equation}\label{eq:liouvillealgebra}
 \mathcal B=\C\oplus\HH_{>0}(\C),\qquad
 \mathcal B/\HH_{>0}(\C)\cong\C.
\end{equation}
More generally, a bound $|\hatF(z)|\leq Mt^\beta$ with a real $M$ exists if and only if $F=t^\beta(c+R)$ with $R$ of positive support.
\end{theorem}
\begin{proof}
Suppose (i). If $F$ had a negative least exponent, choose an ordinary point at which its leading coefficient is nonzero. The value there would be infinite, contradicting the real bound. Thus $F\in\HH_{\geq0}(\C)$. Let $f_0$ be its exponent-zero coefficient, possibly zero. Taking standard parts of the bound at ordinary points gives $|f_0(z)|\leq M$ throughout $\C$. The ordinary Liouville theorem makes $f_0$ a constant $c$. All remaining support is positive, proving (ii).

Conversely, $\widehat R(z)$ is infinitesimal at every finite $z$, so $|\hatF(z)|<|c|+1$ for all such $z$. This proves a single real bound, despite the possible growth of individual higher coefficient functions. The quotient statement follows, and multiplying by $t^{-\beta}$ proves the final assertion.
\end{proof}

\begin{example}[Why the infinitesimal tail cannot be removed]
The function $F(z)=tz$ is nonconstant in $\HH(\C)$, yet $|\hatF(z)|<1$ for every finite surcomplex $z$. More generally, $t f(z)$ is real-bounded on the finite plane for every entire complex function $f$, even when $f$ is unbounded or has arbitrarily rapid classical growth. It has no local maximum unless it is constant.
\end{example}

\begin{corollary}[Full coefficientwise Liouville theorem]
If every coefficient $f_\gamma$ of $F\in\HH(\C)$ is bounded on $\C$, with a bound allowed to depend on $\gamma$, then $F$ is a constant element of $\K$.
\end{corollary}
\begin{proof}
Apply the ordinary Liouville theorem to every coefficient. Their well-ordered Hahn sum is a single surcomplex constant.
\end{proof}
The exact distinction is between bounding the whole value and controlling every asymptotic coefficient. A leading scale can mask unlimited behavior at smaller scales.

\subsection{Schwarz's lemma: a surviving form and an explicit failure}
Let $\Disc=\{z\in\C:|z|<1\}$. The subclass $\Disc^\sharp$ excludes the monads of points of the ordinary unit circle. It is strictly smaller than the full surreal-radius disk $\{z\in\K:|z|<1\}$.

\begin{proposition}[Schwarz lemma for canonical lifts]\label{prop:schwarz}
If $f:\Disc\to\Disc$ is ordinary holomorphic and $f(0)=0$, then
\[
 |f^\sharp(z)|\leq|z|\qquad(z\in\Disc^\sharp).
\]
Equality at a nonzero $z$ holds only when $f(z)=uz$ for a complex constant $u$ with $|u|=1$.
\end{proposition}
\begin{proof}
The ordinary Schwarz lemma gives a holomorphic $g=f/z$, with its removable value at zero, satisfying $|g|\leq1$. If $g$ is not a constant of modulus one, the ordinary maximum-modulus principle gives $|g(c)|<1$ at every $c\in\Disc$. Infinitesimal Taylor evaluation preserves this strict inequality because it has a positive real margin at each fixed $c$. Hence $|g^\sharp(c+\epsilon)|<1$ and $f^\sharp(z)=zg^\sharp(z)$ gives strict inequality for $z\ne0$. The unimodular constant case gives equality.
\end{proof}
For general Hahn-supported disk maps the corresponding assertion is false:
\[
 F(z)=(1+t)z
\]
maps $\Disc^\sharp$ into itself and fixes zero, since its standard part is the identity, but $|F(z)|=(1+t)|z|>|z|$ for every nonzero $z$. In particular its derivative at zero has modulus greater than one. This example identifies both the role of canonical lifting and the importance of distinguishing a thickened ordinary disk from a full surreal-radius disk.

\section{Normal families and harmonic consequences}\label{sec:normal}
\subsection{A coefficientwise Montel theorem}
\begin{definition}
For a fixed well-ordered support set $S$, coefficientwise compact-open convergence means
\[
 F_n=\sum_{\gamma\in S}f_{n,\gamma}t^\gamma\longrightarrow
 F=\sum_{\gamma\in S}f_\gamma t^\gamma
\]
when $f_{n,\gamma}\to f_\gamma$ uniformly on every ordinary compact subset of $U$, separately for each $\gamma\in S$.
\end{definition}
This is a topology on coefficient data, not a claim of pointwise all-radii convergence of evaluated values.

\begin{theorem}[Countable-support Montel theorem]\label{thm:montel}
Let $(F_n)$ have a common countable well-ordered support $S$. Suppose that for each $\gamma\in S$ and each compact $K\Subset U$,
\[
 \sup_n\sup_{z\in K}|f_{n,\gamma}(z)|<\infty.
\]
Then a subsequence converges coefficientwise compact-open to an element of $\HH(U)$. Every fixed coefficientwise derivative converges in the same sense.
\end{theorem}
\begin{proof}
Enumerate $S$. The ordinary Montel theorem gives a locally uniformly convergent subsequence for the first coefficient. Refine for the second coefficient, and continue. The diagonal subsequence converges for every coefficient. Its limit coefficients are holomorphic and retain support in the well-ordered set $S$. Ordinary local uniform convergence of derivatives, obtained from Cauchy's formula, proves the last assertion.
\end{proof}
Countability here is not cosmetic: the diagonal argument does not automatically give a subsequence controlling uncountably many independent coefficients. More generally one can use product-compactness arguments for suitable coefficient spaces and nets, but those nets still concern coefficient topology, not the topology of Theorem~\ref{thm:topology}.

\subsection{Harmonic functions and conjugates}
For an evaluated Hahn-supported holomorphic function $F=u+iv$, the Cauchy--Riemann equations and commuting second derivatives give
\[
 u_{xx}+u_{yy}=0,\qquad v_{xx}+v_{yy}=0.
\]
Here the derivatives are genuine all-radii partial derivatives, by Theorem~\ref{thm:derivative}. Thus real and imaginary parts are harmonic in this differential sense.

Conversely, suppose a common-support Hahn family has real-valued harmonic coefficients on a simply connected ordinary complex domain. Choose the harmonic conjugate of each coefficient with a fixed normalization at one point. The same Hahn support gives a holomorphic family whose real part is the original one. This is a coefficientwise harmonic-conjugate theorem. It is not a claim that every arbitrary class function satisfying a Laplace equation has such a common-support representation.

\section{Other scales and the existing global exponential}\label{sec:exp}
\subsection{Affine charts throughout \texorpdfstring{$\K$}{K}}
Let $a\in\K$ and $\lambda\in\K^\times$. An affine macrochart has domain
\[
 U^\sharp_{a,\lambda}=a+\lambda U^\sharp.
\]
A Hahn-supported function in this chart is
\[
 \mathcal F(a+\lambda w)=\hatF(w),\qquad F\in\HH(U).
\]
Every surcomplex point admits such charts. The derivative and the specified affine contour integral satisfy
\begin{align*}
 \mathcal F'(a+\lambda w)&=\lambda^{-1}\widehat{DF}(w),\\
 \int_{a+\lambda\Gamma}^{\htop}\mathcal F(z)\dd z
 &:=\lambda\int_\Gamma^{\htop}F(w)\dd w.
\end{align*}
These formulas follow from the chain rule and the definition of change of variable. Standard complex affine reparametrizations agree with ordinary coefficientwise substitution. In any additional chart change, the pullback must first be shown to possess an admissible common-support expansion; unrestricted chart-independence is not postulated.

Cauchy's formula, preparation, residues, and root counts therefore apply equally to infinite centers and infinitesimal or infinite scales whenever the displayed chart hypothesis holds. A root monad becomes $a+\lambda(c+\mm)$. Rational functions with arbitrary surcomplex coefficients fit these charts as fractions of finite polynomials.

\subsection{Canonical surcomplex exponentiation}
The following is an existing result, not a construction originating in this article.

\begin{theorem}[Ehrlich--Kaplan, quoted]\label{thm:EK}
There is a canonical surjective group homomorphism
\[
 \operatorname{Exp}:(\K,+)\longrightarrow(\K^\times,\cdot)
\]
extending surreal real exponentiation and finite complex analytic exponentiation, whose kernel is $2\pi i\Oz$. Here $\Oz$ is the omnific integer part of $\No$.
\end{theorem}
See \cite[Proposition 11.2 and Theorem 11.1]{EK}. In normal-form terms, omnific integers have no infinitesimal part, their exponent-zero coefficient is an ordinary integer, and their remaining real normal-form terms are purely infinite. The selection of this integer part is part of the canonical construction; the theorem is not a uniqueness assertion among every conceivable complex exponential extension.

\begin{proposition}[Global differential calculus of $\operatorname{Exp}$]\label{prop:expderivative}
The canonical exponential is microanalytic at every surcomplex point and satisfies
\[
 \operatorname{Exp}'(z)=\operatorname{Exp}(z).
\]
It has a microanalytic logarithm inverse near every nonzero value. On a centered unit-scale chart,
\[
 \operatorname{Exp}(a+w)=\operatorname{Exp}(a)\,(e^w)^\sharp
 \qquad(w\in\Fin),
\]
so the chart expression belongs to $\HH(\C)$.
\end{proposition}
\begin{proof}
For infinitesimal $h$, the established finite analytic extension gives
\[
 \operatorname{Exp}(h)=\Hsum_{n\geq0}\frac{h^n}{n!}.
\]
The homomorphism identity then gives the exact local expansion
\[
 \operatorname{Exp}(z+h)=\operatorname{Exp}(z)
      \Hsum_{n\geq0}\frac{h^n}{n!}.
\]
Differentiation and local invertibility follow from the microscopic results. The finite-chart formula is the same homomorphism identity with finite $w$.
\end{proof}
For a chosen $z_0$ with $\operatorname{Exp}(z_0)=w_0\ne0$, the local logarithm is explicitly
\[
 \operatorname{Log}_{z_0}(w_0(1+\epsilon))
 =z_0+\Hsum_{n\geq1}\frac{(-1)^{n+1}}n\epsilon^n.
\]
Different inverse choices can differ by elements of $2\pi i\Oz$, not only by ordinary integer multiples of $2\pi i$.

Defining
\[
 \operatorname{Sin}z=\frac{\operatorname{Exp}(iz)-\operatorname{Exp}(-iz)}{2i},
 \qquad
 \operatorname{Cos}z=\frac{\operatorname{Exp}(iz)+\operatorname{Exp}(-iz)}2
\]
gives global microanalytic functions. Algebra and differentiation prove the addition formulas,
$\operatorname{Sin}^2z+\operatorname{Cos}^2z=1$,
$\operatorname{Sin}'=\operatorname{Cos}$, and
$\operatorname{Cos}'=-\operatorname{Sin}$. These are consequences of the existing exponential structure and the local calculus, not evidence for an unrestricted global Cauchy theory.

\subsection{A concrete boundary of the common-support theory}
\begin{proposition}\label{prop:movingscale}
The globally microanalytic function $z\mapsto\operatorname{Exp}(\omega z)$ does not belong to $\HH(U)$ on any ordinary connected open domain $U$ containing a nonempty real interval.
\end{proposition}
\begin{proof}
Let $a=\val(\operatorname{Exp}(\omega))<0$, since the real exponential of $\omega$ exceeds every positive real number. For every rational $q$, the homomorphism identity implies
\[
 \val(\operatorname{Exp}(\omega q))=qa.
\]
Suppose a representation $F\in\HH(U)$ existed and let $\beta=\val(F)$, with nonzero holomorphic leading coefficient $f_\beta$. At each ordinary point where $f_\beta$ is nonzero, the value has valuation exactly $\beta$. Among the rational points of the real interval, at most one has $qa=\beta$. Thus $f_\beta$ vanishes at all the other rational points. They accumulate inside $U$, contradicting the complex identity theorem.
\end{proof}
This is not a defect in the proofs; it delineates the selected category. The function has a moving leading scale as its ordinary argument varies. A theory including such functions on a fixed macroscopic chart needs more than a common well-ordered support of fixed monomials, for example a controlled transserial extension. The microscopic and macroscopic levels genuinely differ.

\section{Worked examples and exact verification}\label{sec:examples}
\subsection{Two infinitesimal zeros of a transcendental perturbation}
Consider
\begin{equation}\label{eq:exampleF}
 F(z)=z^2+t e^z\in\HH(\C).
\end{equation}
Its leading coefficient is $z^2$. Therefore $\hatF$ has exactly two zeros in $\Fin$, both in the monad of zero and counted with multiplicity. They are simple: if $F(z)=0$ and $F'(z)=0$, then
\[
 t(e^z)^\sharp=-z^2,\qquad 2z-z^2=0.
\]
A root is infinitesimal and nonzero, so it cannot satisfy $z(2-z)=0$.

Set $s=t^{1/2}$. The equation can be written
\[
 z=A(e^{z/2})^\sharp,\qquad A=\pm is.
\]
Formal Lagrange inversion, or recursive substitution, gives
\begin{equation}\label{eq:rootseries}
 z_\pm=\Hsum_{n\geq1}
       \frac{n^{n-1}}{2^{n-1}n!}(\pm is)^n.
\end{equation}
For completeness, Lagrange inversion in this instance follows by the formal residue change of variables $A=ze^{-z/2}$: the coefficient of $A^n$ in $z(A)$ equals
\[
 \frac1n[X^{n-1}]e^{nX/2}
 =\frac{n^{n-1}}{2^{n-1}n!}.
\]
All coefficients are complex and $s$ is infinitesimal, so the series are Hahn summable. Their first terms are
\begin{align*}
 z_+&=is-\frac{s^2}{2}-\frac{3is^3}{8}
       +\frac{s^4}{3}+\frac{125is^5}{384}-\frac{27s^6}{80}
       +O_{\mathrm{formal}}(s^7),\\
 z_-&=-is-\frac{s^2}{2}+\frac{3is^3}{8}
       +\frac{s^4}{3}-\frac{125is^5}{384}-\frac{27s^6}{80}
       +O_{\mathrm{formal}}(s^7).
\end{align*}
The notation $O_{\mathrm{formal}}(s^7)$ here denotes a formal tail with powers at least seven, not a sequential convergence assertion.

Preparation starts with $E=te^z$ and $m=2$. At first order,
\[
 p^{[1]}=t(1+z),\qquad
 v^{[1]}=t\frac{e^z-1-z}{z^2}.
\]
At second order,
\[
 p^{[2]}=-t^2\left(\frac12+\frac23z\right).
\]
Thus the monic polynomial containing the entire finite zero cluster begins
\[
 P(z)=z^2+t(1+z)-t^2\left(\frac12+\frac23z\right)
       +O_{\mathrm{formal}}(t^3).
\]
The sum of the two root expansions gives the negative linear coefficient, and their product gives the constant coefficient, providing an independent check of preparation.

For every positive real radius $r$,
\[
 \frac1{2\pi i}\int_{|z|=r}^{\htop}
       \frac{2z+te^z}{z^2+te^z}\dd z=2.
\]
This counts the two actual infinitesimal roots, not merely the double zero of an unperturbed limiting polynomial. The finite-plane statement does not exclude other roots in a separate global extension at infinite points.

\subsection{Two genuinely separated surreal scales}
Let $\eta=t^\omega$. Then $\eta\prec t^n$ for every positive integer $n$. Consider
\[
 P(z)=z^2-\eta z-t.
\]
The two roots are exactly
\[
 z_\pm=\frac{\eta\pm\sqrt{\eta^2+4t}}2
 =\pm t^{1/2}+\frac\eta2
   \pm\frac{\eta^2}{8t^{1/2}}
   \mp\frac{\eta^4}{128t^{3/2}}+\cdots.
\]
The binomial expansion is legitimate because $\eta^2/t\in\mm$. The support can involve exponents $\omega$, $2\omega-1/2$, $4\omega-3/2$, and further values not represented by a fixed ordinary power series in $t$ with integer exponents. Nonetheless the preparation and zero-counting theorems apply without modification.

\subsection{Cauchy extraction at a nonstandard target}
Take $F(z)=e^z+t/(2-z)$, a standard unit circle $\Gamma$, and $a=1/3+\epsilon$ with $\epsilon\in\mm$. Then
\[
 \frac1{2\pi i}\int_\Gamma^{\htop}
    \frac{e^z+t/(2-z)}{z-a}\dd z
 =e^{1/3}\Hsum_{n\geq0}\frac{\epsilon^n}{n!}
  +\frac{t}{5/3-\epsilon}.
\]
The circle, coefficients, pole displacement, and returned value belong to different levels of the construction. The formula is nevertheless exact because the joint support is well ordered and locally finite.

\subsection{Reproducibility}
The accompanying \texttt{verify\_examples.py} uses exact rational and symbolic arithmetic to check the Lagrange root series in \eqref{eq:rootseries}, its substitution into \eqref{eq:exampleF} to a specified finite order, the first preparation coefficients, a multiscale binomial identity, and a moving-pole Cauchy extraction. The output is included as \texttt{verification.txt}. These finite checks are diagnostics for examples. The proofs of summability, factorization, and zero counting do not depend on numerical sampling or on a computer algebra system.

\section{What has been achieved, and what would extend it}\label{sec:conclusion}
The theory separates three structures that are easily confused. The class field gives exact algebra and arbitrarily small scales. Hahn summation makes infinitesimal formal substitutions meaningful. A common ordinary complex domain links coefficient functions across otherwise disconnected monads. Contour integration is then defined coefficientwise rather than smuggled in through an unsuitable topology.

Within this category, the central implications have been proved:
\[
 \begin{gathered}
 \text{common-support holomorphy}
 \Longrightarrow\text{Taylor evaluation and differentiation},\\
 \text{positive perturbation}
 \Longrightarrow\text{preparation}
 \Longrightarrow\text{monadic zero conservation},\\
 \text{division and residues}
 \Longrightarrow\text{argument principle and Rouch\'e stability}.
 \end{gathered}
\]
The open mapping and maximum-modulus principles follow from the local algebra together with monadic zero conservation. Liouville's theorem survives exactly after passage to the residue field, or under bounds on every coefficient. A coefficientwise Montel theorem is available under a countable common-support condition. The canonical global exponential has the expected local differential calculus, but its infinite-scale compositions need not fit the macroscopic support condition.

Several extensions are now concrete mathematical projects rather than vague demands for ``complex analysis on all surreals.'' One can seek a larger category admitting moving monomial scales, while preserving enough summability for contour substitution. One can study compatibility of coefficientwise contour integrals with transserial extension and integration constructions. One can develop several-variable preparation and divisor theory with explicit support control. Finally, one can investigate analytic continuation along specified families of overlapping admissible affine charts; ordinary path connectedness in the all-radii topology is not a substitute for that definition.

None of these prospective extensions is assumed in the present theorems. The result is a rigorous, useful surcomplex analysis with explicit boundaries, rather than an unsupported assertion that the entire classical theory transfers unchanged.

\appendix
\section{A hypothesis and dependency audit}\label{app:audit}
\begin{longtable}{>{\raggedright\arraybackslash}p{0.26\textwidth}>{\raggedright\arraybackslash}p{0.31\textwidth}>{\raggedright\arraybackslash}p{0.34\textwidth}}
\toprule
Statement & Essential input & Not being assumed \\
\midrule\endhead
Normal-form evaluation & Set-sized well-ordered support; finite exponent incidence & Ordinary convergence of partial sums.\\
Every formal series is locally realizable & The full surreal set-bound principle; rescaling & The same conclusion over every fixed set-sized Hahn field.\\
Derivative equals Taylor derivative & Joint positive-word support bound & A countable sequence cofinal among surreal small increments.\\
Weierstrass preparation & Positive support of perturbation; coefficientwise Taylor division & Rank-one completeness or cofinality of $n\alpha$.\\
Zero-cluster theorem & Preparation; algebraic closedness of $\K$ & A topological winding number for an all-radii continuous circle.\\
Cauchy formula & Ordinary contour; common coefficient support; target with interior standard part & A contour through boundary monads or arbitrary class paths.\\
Residue theorem & Fraction of supported holomorphic functions; leading denominator nonzero on boundary & Uniform pole order for all expanded coefficients.\\
Maximum modulus & Nonconstant microgerms; global rigidity from the zero-cluster theorem & Compactness of surcomplex closed disks.\\
Liouville replacement & A real bound at every finite point & Bounds on every subleading coefficient.\\
Montel subsequence & Countable common support; coefficientwise local bounds & Convergence in the all-radii value topology.\\
Global exponential & The existing canonical exponential and its finite analytic restriction & A unique exponential independent of integer-part choices.\\
\bottomrule
\end{longtable}

\section{Conventions for using the archive}
The main source is \texttt{surcomplex\_analysis.tex}; the bibliography is embedded so that no separate bibliography processor is needed. Build the article by running \texttt{pdflatex} twice, or by running \texttt{latexmk -pdf surcomplex\_analysis.tex}. The included script uses Python 3 and SymPy. It does not implement the proper class of surreal numbers. It verifies finite formal identities used in the examples.

The mathematical status of the article is a research-style development with explicit proofs and source attribution, not an independently refereed publication. Its claims should be checked on their stated hypotheses. In particular, the separation between formal summation and all-radii convergence, and between ordinary domains and their monadic thickenings, is integral to every application.

\begin{thebibliography}{99}
\bibitem{Conway}
J. H. Conway, \emph{On Numbers and Games}, second edition, A K Peters, 2001; first edition, Academic Press, 1976.

\bibitem{Gonshor}
H. Gonshor, \emph{An Introduction to the Theory of Surreal Numbers}, London Mathematical Society Lecture Note Series 110, Cambridge University Press, 1986.

\bibitem{Alling}
N. L. Alling, \emph{Foundations of Analysis over Surreal Number Fields}, North-Holland Mathematics Studies 141, North-Holland, 1987. In particular, the treatment of formal power series and Neumann's lemma in Chapter 7.

\bibitem{vdDE}
L. van den Dries and P. Ehrlich, Fields of surreal numbers and exponentiation, \emph{Fundamenta Mathematicae} \textbf{167} (2001), 173--188. DOI: \href{https://doi.org/10.4064/fm167-2-3}{10.4064/fm167-2-3}.

\bibitem{vdDEerr}
L. van den Dries and P. Ehrlich, Erratum to ``Fields of surreal numbers and exponentiation'', \emph{Fundamenta Mathematicae} \textbf{168} (2001), 295--297.

\bibitem{RSS}
S. Rubinstein-Salzedo and A. Swaminathan, Analysis on surreal numbers, \emph{Journal of Logic and Analysis} \textbf{6} (2014), paper 5, 1--39. \href{https://arxiv.org/abs/1307.7392}{arXiv:1307.7392}.

\bibitem{EK}
P. Ehrlich and E. Kaplan, Surreal ordered exponential fields, \emph{Journal of Symbolic Logic} \textbf{86} (2021), 1066--1115. \href{https://arxiv.org/abs/2002.07739}{arXiv:2002.07739}; especially Section 11 on surcomplex exponentiation.

\bibitem{CE}
O. Costin and P. Ehrlich, Integration on the surreals, \emph{Advances in Mathematics} \textbf{452} (2024), 109823. \href{https://arxiv.org/abs/2208.14331}{arXiv:2208.14331}.

\bibitem{MS}
J. M\"uller and A. Strohmaier, The theory of Hahn-meromorphic functions, a holomorphic Fredholm theorem, and its applications, \emph{Analysis \& PDE} \textbf{7} (2014), 745--770. DOI: \href{https://doi.org/10.2140/apde.2014.7.745}{10.2140/apde.2014.7.745}.

\bibitem{Marker}
D. Marker, \emph{Model Theory of Valued Fields}, lecture notes, University of Illinois at Chicago, 2018, Sections 1--2, especially Lemma 2.36. Available from the author's course page:
\url{https://homepages.math.uic.edu/~marker/math512-f18/valued_fields_1-2.pdf}.

\bibitem{Ahlfors}
L. V. Ahlfors, \emph{Complex Analysis}, third edition, McGraw--Hill, 1979. Used for the ordinary complex Cauchy, identity, Liouville, and normal-family theorems underlying the coefficient arguments.
\end{thebibliography}
\end{document}
